%!TEX root = ../thesis.tex
\chapter{Introduction}
\label{chap:introduction}

\noindent

% ============================================================================
\section{Motivation}
\label{sec:motivation}

End-to-end autonomous driving learns a single visual-to-motion stack, avoiding hand-engineered interfaces between perception, prediction, and planning. Under the standard pooled evaluation protocol, the paradigm appears mature: recent planners evaluated on NAVSIM report strong aggregate PDMS and match or exceed classical modular baselines on geographically mixed test sets~\cite{dauner-navsim-2024}. The protocol itself is the caveat. A model that interpolates well across pooled data from Boston, Pittsburgh, Las Vegas, and Singapore is not necessarily robust to deployment in a city it has not seen.

Once evaluation is re-framed as zero-shot geographic transfer, the picture changes. Naeinian et al.~\cite{naeinian-cross-city-2026} report order-of-magnitude open-loop degradation for supervised Boston-to-Singapore transfer, while closed-loop NAVSIM results show smaller out-of-distribution drops for self-supervised features. Chapter~\ref{chap:experiments} reports the exact matrices and ratios; the point here is conceptual. These failure modes are safety-critical: an autonomous vehicle trained in one city must remain competent when moved to another without expensive re-collection and re-labelling of local data.

This thesis studies how visual-representation choice mediates these failures. The central claim is bounded rather than absolute:

\begin{quote}
\textbf{Under explicit city-disjoint evaluation, end-to-end driving failures are strongly mediated by the visual representation. In the evaluated settings, frozen self-supervised backbones transfer more reliably across cities than supervised baselines, while planner choice appears secondary to representation choice.}
\end{quote}

\noindent
The scope of this claim is deliberately tied to the evaluated settings (\reftab{tab:intro-rqs}). The thesis does not claim that planner architecture is irrelevant, nor that self-supervised features are city-invariant by construction; it defends a narrower and more defensible version of the representation-over-architecture argument on the evidence currently completed.

% ============================================================================
\section{Scope and Framing}
\label{sec:project-evolution}

The thesis began with a narrower question: whether a frozen self-supervised encoder paired with a small MLP waypoint regressor could already improve NAVSIM planning under mixed-city evaluation~\cite{hamza-sanchez-navsim-report}. Drive-JEPA~\cite{wang-drive-jepa} subsequently showed that V-JEPA-style video pre-training is already highly competitive on pooled NAVSIM, making mixed-city PDMS an insufficient discriminator for a representation-focused thesis. The scope shifted from \emph{whether} SSL helps mixed-city planning to \emph{where} it matters most and \emph{why}, with geographic transfer identified as the regime in which representation choice carries the most diagnostic weight.

% ============================================================================
\section{Problem Statement}
\label{sec:problem}

The subject of study is geographic domain shift in end-to-end planning. The cities used in this thesis differ in road topology, lane geometry, infrastructure layout, traffic density, and driving convention (right-hand traffic in Boston, Pittsburgh, and Las Vegas; left-hand traffic in Singapore). These structural differences create a harder transfer setting than mixed-city validation, because city-specific shortcuts learned during training do not transfer~\cite{naeinian-cross-city-2026}.

The thesis isolates the \textbf{representation axis} as its primary axis: holding the planner, training recipe, and evaluation protocol fixed, we vary the visual backbone across supervised and self-supervised initialisations and measure cross-city behaviour. A secondary \textbf{planner-robustness test} asks whether the representation signal observed under regression-based planners (LAW, TransFuser, Latent TransFuser) also survives under a hostile-to-backbone generative head (DiffusionDrive). The mixed-city portion of this test is completed and reported in Chapter~\ref{chap:experiments}; the corresponding city-disjoint DiffusionDrive evaluation is specified as active validation with pre-registered success criteria, not as part of the completed argument.

% ============================================================================
\section{Research Questions}
\label{sec:research-questions}

\reftab{tab:intro-rqs} lists the questions the completed thesis answers, together with the section of Chapter~\ref{chap:experiments} that provides the evidence. Follow-on questions whose evaluation is active are identified separately in \refsec{sec:contributions} and are not included in the formal research-question table.

\begin{table}[ht]
\centering
\caption{Completed research questions and the chapter section that addresses each.}
\label{tab:intro-rqs}
\small
\setlength{\tabcolsep}{6pt}
\renewcommand{\arraystretch}{1.25}
\begin{tabular}{@{}>{\raggedright\arraybackslash}p{1.1cm} >{\raggedright\arraybackslash}p{7.5cm} >{\raggedright\arraybackslash}p{4.2cm}@{}}
\toprule
\textbf{ID} & \textbf{Question} & \textbf{Addressed in} \\
\midrule
RQ1 & Which self-supervised pre-training objective (I-JEPA, DINOv2, MAE) transfers best across cities relative to supervised ImageNet features, under a fixed planner? & Cross-city SSL benchmark \\
RQ2 & Which NAVSIM city (Boston, Las Vegas, Pittsburgh, Singapore) is the most reliable single-city training source for zero-shot transfer? & Cross-city SSL benchmark \\
RQ3 & Do frozen self-supervised features degrade more gracefully than supervised baselines under city shift, and which mechanistic probes (probe accuracy, CKA, MMD, effective rank) track the ranking? & Cross-city benchmark + mechanistic analysis \\
RQ4 & Does the representation-centric pattern observed under regression planners reproduce under a mixed-city generative planner (DiffusionDrive), or is it regression-head specific? & Generative-planning study (mixed-city only) \\
\bottomrule
\end{tabular}
\end{table}

% ============================================================================
\section{Contributions}
\label{sec:contributions}

The thesis contributions are separated into completed contributions and active validation. Prior-work boundaries are stated explicitly.

\paragraph{Completed contributions.}
\begin{enumerate}
    \item \textbf{NAVSIM City Splits}. A per-city train/test partitioning of NAVSIM across Boston, Pittsburgh, Las Vegas, and Singapore. The partition and its open/closed-loop evaluation protocol were introduced in our earlier cross-city study~\cite{naeinian-cross-city-2026}; the thesis integrates that protocol into a broader representation-centred investigation and adds a NAVSIM v1.1 re-evaluation that hardens the numbers against devkit-version confounding.
    \item \textbf{Cross-city representation benchmark}. A controlled evaluation of self-supervised backbones (I-JEPA, DINOv2, MAE) against supervised ResNet-34 and Swin-T baselines across LAW on nuScenes and TransFuser / Latent TransFuser on NAVSIM. Under the evaluated settings, frozen self-supervised backbones reduce cross-city degradation relative to supervised baselines; the exact ratios and matrices are reported in Chapter~\ref{chap:experiments}.
    \item \textbf{DiffusionDrive integration (mixed-city supporting breadth)}. Self-supervised backbones (I-JEPA, MAE) were integrated into the DiffusionDrive generative head under the same feature adapter used in the cross-city benchmark, including a camera-only variant. The completed mixed-city results show that the representation signal survives a planner swap under pooled evaluation, at a small in-distribution cost relative to the supervised C+L baseline. This is reported as supporting breadth, not as completed cross-city evidence.
    \item \textbf{Mechanistic analysis}. A study of the trained checkpoint roster using linear probing for city identity, centred kernel alignment, maximum mean discrepancy, and feature-covariance effective rank. Within the evaluated settings, the linear-probe ordering mirrors the out-of-distribution PDMS drop ordering across matched backbone families, and fine-tuning compresses feature-covariance rank for every self-supervised family; objective-specific nuance is reported without claiming a universal mechanism.
\end{enumerate}

\paragraph{Active validation (not claimed as completed evidence).}
\begin{itemize}
    \item \textbf{Cross-city generative planning}. The DiffusionDrive cross-city matrix is specified in Chapter~\ref{chap:methodology} with pre-registered success criteria and is scheduled for execution. The thesis does not use its outcome as load-bearing evidence.
    \item \textbf{I-JEPA camera-only DiffusionDrive}. The mixed-city camera-only I-JEPA run is specified with reading bands in Chapter~\ref{chap:experiments}. Its outcome is reported if and when the run completes; it is not part of the defended contribution set.
    \item \textbf{RL post-training and cross-city transfer}. Whether GRPO / DPO gains on DiffusionDrive variants transfer across cities is flagged as a follow-on question only.
\end{itemize}

% ============================================================================
\section{Thesis Organization}
\label{sec:organization}

The remainder of the thesis is organised as follows.

\begin{itemize}
    \item \textbf{Chapter~\ref{chap:background}} covers background on end-to-end autonomous driving, the NAVSIM benchmark, the self-supervised pre-training objectives evaluated (I-JEPA, DINOv2, MAE), and the planning architectures used (LAW, TransFuser, Latent TransFuser, DiffusionDrive).
    \item \textbf{Chapter~\ref{chap:relatedwork}} positions the thesis against related work in self-supervised learning for driving, cross-domain generalisation, and diffusion-based planning, including how Drive-JEPA altered the in-distribution performance landscape during the course of the work.
    \item \textbf{Chapter~\ref{chap:methodology}} presents the methodology: NAVSIM City Splits construction, the shared backbone-adapter interface, the cross-city evaluation protocol, the DiffusionDrive integration framed as a hostile-to-backbone planner test, and the active-validation specification for the cross-city generative study.
    \item \textbf{Chapter~\ref{chap:experiments}} reports the completed experimental evidence: the exploratory lightweight-head study, the cross-city SSL benchmark, the mixed-city DiffusionDrive integration, and the mechanistic analysis.
    \item \textbf{Chapter~\ref{chap:conclusions}} summarises the completed contributions, states the claim boundaries, and identifies the open validation and follow-on questions.
\end{itemize}
